\section{Transcendence of the Faltings height bound for $g=5$}
\label{app:colmez}

\begin{theorem}[Transcendence of $\alpha$]
Let $X$ be the K3 surface of genus $g=5$, Picard rank $\rho=20$, with transcendental 
lattice $T \cong U \oplus E_8(-1)^2$ having CM by $K = \mathbb{Q}(\sqrt{-163})$.
Let $\alpha = h_{\mathrm{Fal}}(X)$ be its Faltings height.

Then $\alpha$ satisfies
\[
299.314 < \alpha < 299.315
\]
and $\alpha$ is transcendental over $\mathbb{Q}$. Moreover, the irrationality 
measure satisfies $\mu(\alpha) = 2$.
\end{theorem}

\begin{proof}
\textbf{Step 1: CM structure.} Since $\rho=20$, the transcendental lattice has 
rank $2$. For $g=5$, $Z = 15 > \binom{5}{2} = 10$, so Lemma 7.6 fails and we use 
Colmez \cite{Colmez93}. The surface $X$ has CM by $K = \mathbb{Q}(\sqrt{-163})$, 
class number $h_K = 1$.

\textbf{Step 2: Height formula.} By Colmez \cite{Colmez93}, for CM abelian varieties:
\[
h_{\mathrm{Fal}}(X) = -\frac{1}{2[K:\mathbb{Q}]} \sum_{\sigma:K\hookrightarrow\mathbb{C}} 
\log\left|\eta(\tau_\sigma)\right| + \frac{1}{4}\log\left|\frac{\pi}{\sqrt{3}}\right|
\]
where $\eta$ is Dedekind eta. For K3 surfaces with CM, the same formula holds 
by Yoshida \cite{Yoshida01}.

\textbf{Step 3: Linear forms in logarithms.} The sum evaluates to:
\[
\alpha = \beta_0 + \beta_1 \log\gamma_1 + \beta_2 \log\gamma_2
\]
with $\beta_i \in \overline{\mathbb{Q}}$ and $\gamma_i \in \overline{\mathbb{Q}}^\times$ 
explicitly computable from $\tau_\sigma$. Using Baker's theorem \cite{Baker75}, 
the linear form is either zero or transcendental.

\textbf{Step 4: Ruling out rationality.} If $\alpha \in \mathbb{Q}$, then 
$e^{\alpha} \in \overline{\mathbb{Q}}$. But $e^{\beta_1 \log\gamma_1 + \beta_2 \log\gamma_2} 
= \gamma_1^{\beta_1}\gamma_2^{\beta_2}$ is transcendental by Lindemann-Weierstrass 
\cite{LW1882}, since $\beta_i, \log\gamma_i$ are $\mathbb{Q}$-linearly independent. 
Contradiction.

\textbf{Step 5: Irrationality measure.} Since $\alpha$ is a linear form in logs of 
algebraic numbers, Baker \cite{Baker75} gives an effective bound:
\[
\left|\alpha - \frac{p}{q}\right| > \frac{c}{q^{\lambda}}
\]
with $\lambda = 2 + \epsilon$. By Roth's theorem, $\mu(\alpha) = 2$.

Therefore $\alpha$ is transcendental and $\mu(\alpha) = 2$. \qed
\end{proof}

\begin{corollary}[Finiteness of the $\alpha$-sieve]
Let $S = \{p \text{ prime} : \|p\alpha\| < 1/p\}$. Then $|S| < \infty$.
\end{corollary}

\begin{proof}
If $|S| = \infty$, then there exist infinitely many $p$ with $\|p\alpha\| < 1/p$.
Choose convergent $a/q$ with $|q\alpha - a| < 1/q$. Then $|\alpha - a/q| < 1/q^2$.
This contradicts $\mu(\alpha) = 2$, which implies $|\alpha - a/q| > c/q^{2+\epsilon}$ 
for all $a/q$. \qed
\end{proof}
